\documentclass[11pt]{article}
\usepackage[T1]{fontenc}
\usepackage{lmodern}
\usepackage[margin=0.92in]{geometry}
\usepackage{amsmath,amssymb,amsthm,booktabs,microtype}
\usepackage[colorlinks=true,linkcolor=blue,citecolor=blue,urlcolor=blue]{hyperref}
\hypersetup{pdftitle={Scalar-level initial geometry and robust shell-volume perturbations},pdfauthor={Ricardo Maldonado},pdfsubject={Working manuscript for author and specialist review}}
\pdfinfoomitdate=1
\pdftrailerid{}
\pdfsuppressptexinfo=15
\newcommand{\eps}{\varepsilon}
\newcommand{\dd}{\mathrm d}
\newcommand{\Lie}{\mathcal L}
\newcommand{\av}[1]{\langle #1\rangle}
\newcommand{\norm}[1]{\lVert #1\rVert}
\newcommand{\cV}{\mathcal V}
\newtheorem{theorem}{Theorem}
\newtheorem{lemma}{Lemma}
\newtheorem{proposition}{Proposition}
\title{Scalar-level initial geometry and robust shell-volume perturbations}
\author{Ricardo Maldonado\\\small Independent Researcher}
\date{4 November 2026}

\begin{document}
\maketitle
\begin{center}\small Working manuscript for author and specialist review\end{center}
\begin{abstract}
We describe a geometric, background-calibrated class of homogeneous second
preparations in a conditional five-dimensional Einstein--scalar shell model.
On a physically supplied initial Cauchy section, we compare the three-area
of the continued spherical component of each scalar level with that of a
fixed base preparation. For families with the same first tangent, the
relative difference of second area derivatives is
$D=3J/2-3aw/p$, including the scalar displacement of the level.
In the retained compact metric-mean freedom, with the other canonical
sectors fixed, the full profile determines the preparation by a two-component
linear initial-value problem. On the target's causal initial segment,
exact rational bounds give $N\le R_A/2000$, where $R_A$ bounds $D$ and
its fixed background scalar-label derivative. The specified class
$R_A\le2\times10^{33}$ therefore retains strictly negative proper-volume
and expansion second coefficients at one fixed proper time. Each admitted
exact family consequently has smaller positive volume and slower positive
expansion for sufficiently small nonzero parameter of either sign.
The result retains explicit background, constraint, development, and
uniqueness premises. It supplies neither a physical rule choosing this
class nor a parameter interval uniform over it. The contribution is the
restricted geometric reconstruction and its quantitative transport to these
two model observables, rather than a general gauge-invariance principle.
\end{abstract}

\section{Question, model, and retained premises}

A fixed first perturbation does not fix homogeneous second initial data.
In the present shell model, an admissible compact second preparation can
reverse the signs of both proper-volume and expansion coefficients while
preserving the first packet and the initial shell collar
\cite{preparation,records}. This raises a narrower question: can a specified
class of second preparations be expressed using initial geometry, and can
one certify both signs throughout that class? We answer this conditionally
for the retained metric-mean freedom. The initial physical section and a
background scalar-label calibration remain part of the specification.

The bulk action has normalization
\begin{equation}
 S_{\rm bulk}=\int\left(\tfrac12R-\tfrac12|\dd\phi|_g^2-\cV(\phi)\right)
 \dd\mu_g.
 \label{eq:action}
\end{equation}
The bulk and one-copy outward shell equations are
\begin{align}
 \operatorname{Ric}_{AB}-\phi_A\phi_B-\tfrac23\cV(\phi)g_{AB}&=0,
 &\Box_g\phi-\cV'(\phi)&=0,\label{eq:bulk}\\
 K^{\rm out}_{ab}&=\tfrac16\sigma(\phi)\gamma_{ab},
 &n_{\rm out}\phi&=-\tfrac12\sigma'(\phi).\label{eq:shell}
\end{align}
There are two identical interior copies with their respective outward
normals. Fix
\begin{gather}
 W(\phi)=1-\phi+\phi^3/3,\qquad
 \cV=\tfrac12(W')^2-\tfrac23W^2,\qquad
 \sigma=2W+\delta_\sigma(1+c_\sigma\phi),\nonumber\\
 \delta_\sigma=10^{-3},\qquad
 c_\sigma=96714321425692/161839258930731.
 \label{eq:potential}
\end{gather}
The selected regular near-$\phi=1$ background is
\begin{equation}
 g_0=\dd y^2+\rho(y)^2[-\dd t^2+\cosh^2t\,\kappa],\qquad
 a=\rho_y/\rho>0,\quad p=\phi_{0,y}\ne0,
 \label{eq:background}
\end{equation}
where $\kappa$ is the unit round metric on $S^3$. Its exact selected shell
radius is $\rho_b$. Put $\ell=10^{-6}$, $x=t/\ell$, $r=z/\ell$, and
$\dd y=-\rho\,\dd z$. The packet is $1\le r\le2$; the target is
\begin{equation}
 T=5/4,\qquad \tau_*=\rho_b\ell T.
 \label{eq:target}
\end{equation}
The exact $\rho_b$, rather than a rounded value from an enclosure, specifies
this clock. All quantities use the inherited model and coefficient units.

Every comparison keeps the selected background, first seed and parameter
normalization, full quadratic constraint and evolution sources, canonical
momentum and scalar canonical second sectors, second transverse-traceless
initial pair, and initial shell collar fixed. In particular,
\begin{equation}
 x_\eps=x_0+\eps v_{\rm seed}+\tfrac12\eps^2x_2+o(\eps^2),\qquad
 g_\eps=g_0+\eps q+\tfrac12\eps^2k+o(\eps^2).
 \label{eq:normalization}
\end{equation}
The fixed first seed is a smooth flat quadrupole packet with
$Y=n_1^2-n_2^2$, $\av Y=0$, $\av{Y^2}=1/6$, and
$\Delta_\kappa Y=-8Y$. Its normalization and the base preparation are
identified in Appendix~\ref{app:base}. Its full first-field solution is
common to all families; no quadratic first-field term is discarded in an
absolute observable.

We retain the selected background and its continuous authenticated annulus
bounds; the equivariant weighted constraint isomorphism and nonlinear
implicit-function hypotheses; the conditional fixed-shell development,
subsidiary and boundary uniqueness, and identification of the first
derivative with that first-field solution. For each fixed preparation we
assume $C^2$ dependence on $\eps$ in the retained lower $C^4$ spacetime
topology on a common compact neighborhood, including the required traces.
After a family-dependent shrinkage in $\eps$, an intrinsic normal tube
reaches \eqref{eq:target}. The complete five-mean linear mixed problem
retains its physical Robin and normal-gauge rows, uniqueness, and the
coefficient estimate stated in Section~\ref{sec:transfer}. These are shared
analytic premises of the retained construction \cite{records}; the finite
rational checks below do not establish them. No effective or uniform
nonlinear parameter radius is assumed.

\section{Restricted homogeneous preparation freedom}

On the physical initial section $\Sigma$, the background spatial data are
$\gamma_0=\dd y^2+\rho^2\kappa$ and $\phi_0=\phi_0(y)$. Let $h$ be smooth,
supported in the packet and flat at its endpoints, with
\begin{equation}
 \int_{y_L}^{y_U}\rho^2h\,\dd y=0,\quad
 \chi(y)=\int_{y_L}^{y}\rho^2h\,\dd s,\quad
 h=\rho^{-2}\chi_y,\qquad y_L=y(2),\quad y_U=y(1).
 \label{eq:admissible}
\end{equation}
Thus $\chi$ is smooth, compact, and flat. The complete added second
preparation is the metric mean
\begin{equation}
 \Delta\gamma_2=h\gamma_0+\Lie_{\zeta\partial_y}\gamma_0,
 \quad\Delta\phi_2=0,\qquad
 \zeta=\frac{3(h_y-ah)}{2p^2}+\frac{9\chi}{2\rho^4p^2};
 \label{eq:raw}
\end{equation}
the canonical momentum and both scalar canonical second sectors are fixed.
The background identity $a_y+4a^2=3/\rho^2-2\cV/3$ gives
\begin{align}
 C[h]&=-\tfrac32(h_{yy}+3ah_y)-\cV h,\nonumber\\
 \rho^4C[h]&=-\tfrac32[\rho^4(h_y-ah)]_y-\tfrac92\rho^2h,
 &P_g\Delta\gamma_2&=C[h]+\rho^{-4}(\rho^4p^2\zeta)_y=0.
 \label{eq:constraint}
\end{align}
The moment and flatness conditions make this a compact homogeneous
constraint datum. Adding it to the base preparation leaves the complete
quadratic constraint source unchanged. The retained weighted
implicit-function construction and equivariance give an exact completed
family with the stated second mean for each fixed admissible $h$
\cite{preparation}. This is an application of the premises, not a new
existence theorem for unrestricted data.

The freedom is not the pure second conditioning diffeomorphism. At the
time-symmetric background, a normal generator changes the initial spatial
metric and scalar by multiples of $K_0=0$ and $\Pi_0=0$. Averaging a proposed
tangential generator leaves a radial one; its scalar variation is $pX^y$.
Since the raw scalar difference in \eqref{eq:raw} vanishes and $p\ne0$,
$X^y=0$. Also $\Delta\gamma_2=0$ would require
$h+2\zeta_y=h+2a\zeta=0$, hence $\zeta_y=a\zeta$; compact support then
gives $\zeta=h=0$. Nonzero $h$ is therefore physical in the retained
anchored pure second gauge quotient, with the first representative fixed.

\section{Scalar-level area on the supplied section}

For a fixed background scalar value $\varphi=\phi_0(y)$, let
$S_\eps(\varphi)$ be the component of $\{\phi_\eps=\varphi\}$ in the
selected annular collar that continues the background three-sphere. Define
\begin{equation}
 A_\eps(\varphi)=\operatorname{Vol}_3(S_\eps(\varphi),\gamma_\eps),
 \qquad A_0(\phi_0(y))=2\pi^2\rho(y)^3,
 \qquad
 D_h(y)=\frac{\partial_\eps^2 A_\eps^h(\phi_0(y))|_0
                 -\partial_\eps^2 A_\eps^{\rm base}(\phi_0(y))|_0}
                {A_0(\phi_0(y))}.
 \label{eq:area}
\end{equation}
The scalar label is held fixed in the parameter derivatives. They are
actual second derivatives, with no extra factor $1/2$. The component
qualification excludes unrelated components elsewhere on $\Sigma$.

\begin{lemma}[Regularity and shared-first cancellation]\label{lem:area}
Under the retained $C^2$-into-$C^4$ hypothesis, \eqref{eq:area} exists
and is $C^1$ in the background label. If two preparations have identical
first metric and scalar fields, then
\begin{equation}
 D_h=\av{\tfrac12\operatorname{tr}_{\rho^2\kappa}
                  (\Delta\gamma_{2,\mathrm{tan}})
                    -3a\Delta\phi_2/p}.
 \label{eq:difference}
\end{equation}
Brackets denote the normalized $\kappa$ integral.
\end{lemma}

\begin{proof}
Work in a slightly enlarged compact collar, so endpoint labels are included.
The uniform nonzero bound on $p$ and $C^4$ closeness give unique continued
level graphs $F_\eps(y,\omega)=(u_\eps(y,\omega),\omega)$ satisfying
\begin{equation}
 \phi_\eps(u_\eps(y,\omega),\omega)=\phi_0(y),\qquad u_0=y.
 \label{eq:graph}
\end{equation}
The spatial implicit equations give continuity of $u_\eps$ in $C^4$.
At the graph, its parameter derivatives satisfy
\begin{equation}
 \dot u=-\dot\phi/\phi_y,\qquad
 \ddot u=-\frac{\ddot\phi+2\dot\phi_y\dot u+\phi_{yy}\dot u^2}{\phi_y}.
 \label{eq:graphderivatives}
\end{equation}
Their spatial derivatives through order two use at most four spatial
derivatives of the input and uniformly nonzero denominators. Difference
quotients and the mean value formula give $u_\eps$ $C^2$ into $C^2$.
This allows two spatial orders of loss; differentiability of composition
at unchanged $C^4$ regularity is not needed.

The induced metric, with all fields evaluated at $F_\eps$, is
\begin{equation}
 \sigma_{\eps,AB}=\gamma_{\eps,AB}
 +\gamma_{\eps,yA}\partial_Bu_\eps
 +\gamma_{\eps,yB}\partial_Au_\eps
 +\gamma_{\eps,yy}\partial_Au_\eps\partial_Bu_\eps.
 \label{eq:induced}
\end{equation}
It is $C^2$ into $C^1$ and remains positive definite. Integrating
$\sqrt{\det\sigma_\eps}$ over the compact sphere gives $A_\eps$ $C^2$
into $C^1$. No rotational symmetry of the perturbed data is assumed.

The first level displacements coincide. Subtracting
\eqref{eq:graphderivatives} at $\eps=0$ gives
$\Delta\ddot u=-\Delta\phi_2/p$. The background angular graph slope and
$\gamma_{0,yA}$ vanish, so
\begin{equation}
 \Delta\sigma_{2,AB}=\Delta\gamma_{2,AB}
                 -(\Delta\phi_2/p)\partial_y\gamma_{0,AB}.
 \label{eq:sigmadifference}
\end{equation}
The determinant terms quadratic in the common first variation cancel
after subtraction. Its remaining linear trace gives
\eqref{eq:difference}. The actual nonspherical first-field quadratics
remain in each absolute area variation.
\end{proof}

For an angular metric mean $J$ and scalar mean $w$, the formula becomes
\begin{equation}
 \boxed{D_h=\tfrac32J-3aw/p.}
 \label{eq:invariant}
\end{equation}
In \eqref{eq:raw}, $J=h+2a\zeta$ and $w=0$. After the retained pure
second conditioning gauge, $J=h$ and $w=-p\zeta$. Both give
\begin{equation}
 D_h=\tfrac32(h+2a\zeta),\qquad
 w=\frac{p}{2a}(h-2D_h/3).
 \label{eq:conditioning}
\end{equation}
The scalar pullback in the conditioned representative is essential.
For a general pure second tangential generator $X$, the metric trace in
\eqref{eq:difference} contributes $3aX^y+\operatorname{div}_\kappa X^{\rm tan}$
and the scalar term contributes $-3aX^y$. The sphere integral of the
divergence vanishes. Thus the area is invariant under section
diffeomorphisms transporting both fields, including nonspherical ones.
A spacetime diffeomorphism must also transport the supplied physical
section. Replacing it by a different physical slice changes the
observable; no invariance under that replacement is asserted.

\section{Reconstruction and the geometric class}

Set $Z=\chi/(\ell\rho^3)$, $b=\ell\rho a$, and
$g=\ell\rho p^2/(3a)$. Substitution of \eqref{eq:raw} into
\eqref{eq:conditioning} gives
\begin{equation}
 h_y=(a-p^2/(3a))h-3\chi/\rho^4+(p^2/(3a))(2D_h/3).
 \label{eq:hy}
\end{equation}
Using $\rho_r=-\ell a\rho^2$ and $\chi_r=-\ell\rho^3h$, we obtain
\begin{equation}
 \boxed{\quad Z_r=3bZ-h,\qquad
 h_r=3\ell^2Z+(-b+g)h-\tfrac23gD_h,\qquad
 (Z,h)(1)=(0,0).\quad}
 \label{eq:reconstruction}
\end{equation}
Flatness gives $h(1)=0$ and the exact zero moment gives $\chi(1)=0$.

\begin{proposition}[Restricted injectivity]\label{prop:injectivity}
Within \eqref{eq:admissible}--\eqref{eq:raw}, with the other canonical
second sectors fixed, equal profiles on $[1,T]$ give equal $h$, $\chi$,
and $\zeta$ on that segment. Equal full profiles on $[1,2]$ give equal
full preparations. In particular, $D_h=0$ on the whole packet implies
$h=0$.
\end{proposition}
\begin{proof}
The continuous-coefficient linear initial-value problem
\eqref{eq:reconstruction} has a unique solution by its fundamental matrix
or the integral Gronwall inequality. Its difference for equal profiles
has zero forcing and initial value. The formula for $\zeta$ then gives
the claim.
\end{proof}

A profile zero only on $[1,T]$ need not restrict a preparation outside
that causal segment. This is injectivity for the specified mean freedom,
not for arbitrary Einstein--scalar Cauchy data. It is not surjectivity
onto arbitrary $C^1$ profiles: compactness, endpoint flatness, and the
moment condition constrain the image.

The derivative used to calibrate the class is fixed by the background
scalar label:
\begin{equation}
 D_{h,r}:=\phi_{0,r}\partial_\varphi D_h
        =(-\ell\rho p)\partial_\varphi D_h,
 \qquad \varphi=\phi_0(y(r)).
 \label{eq:labelderivative}
\end{equation}
It is not a derivative in an arbitrary perturbed radial coordinate.
The background calibration can itself be read from
$\rho(\varphi)=(A_0(\varphi)/(2\pi^2))^{1/3}$ and the retained map
$r\mapsto\varphi$. Define
\begin{equation}
 R_A(h)=\max\left\{\sup_{1\le r\le T}|D_h(r)|,
                         \sup_{1\le r\le T}|D_{h,r}(r)|\right\},\qquad
 \mathcal C_A=\{h\text{ satisfying \eqref{eq:admissible}}:
                     R_A(h)\le2\times10^{33}\}.
 \label{eq:class}
\end{equation}
This is a geometric, background-calibrated specification on the given
section. Neither the restriction to these canonical sectors nor the
chosen radius is derived from a physical preparation process. If an
increment is written $\lambda h$, its class norm is
$R_A(\lambda h)=|\lambda|R_A(h)$; the unscaled direction does not control
arbitrary $\lambda$.

\section{An exact rational initial-data bound}\label{sec:rational}

Write $R=R_A(h)$. The retained continuous background boxes imply, everywhere
on the causal initial segment,
\begin{gather}
 129\le\rho,\rho_b\le131,\qquad
 3/25000\le-p\le1/7000,\qquad
 111/1000\le a\le14/125,\nonumber\\
 |p_y|\le1/4000,\qquad |a_y|\le1/10000,\qquad
 1/1200\le-F'_b\le1/1000.
 \label{eq:caps}
\end{gather}
Here $F'_b$ is the scalar-field derivative at the background shell of
$F(\phi)=2\cV(\phi)-\sigma'(\phi)^2/4+\sigma(\phi)^2/3$, the coefficient
in the retained proper-volume difference equation below. These are bounds on continuous
domains, rather than point samples. Their source bytes, the class radius,
and the numerical gates were fixed before the new rational evaluation.

Let $b_+=917/62500000$ and $g_+=131/16317000000000$, obtained directly
from \eqref{eq:caps}. The matrix of \eqref{eq:reconstruction} has
infinity norm at most
\begin{equation}
 \max\{1+3b_+,3\ell^2+b_++g_+\}\le\mu:=20001/20000.
 \label{eq:matrixbound}
\end{equation}
Put $\Delta=T-1=1/4$. A positive Taylor sum through degree six with a
geometric tail gives the exact rational certificate
\begin{equation}
 e^{\mu\Delta}\le
 \sum_{j=0}^6\frac{(\mu\Delta)^j}{j!}
 +\frac{(\mu\Delta)^7/7!}{1-\mu\Delta/8}<\frac43.
 \label{eq:exponential}
\end{equation}
Variation of constants therefore bounds both $|Z|$ and $|h|$ by $HR$,
where
\begin{equation}
 H=\Delta\tfrac43\tfrac23g_+=\frac{131}{73426500000000},\qquad
 H_r=3\ell^2H+(b_++g_+)H+\tfrac23g_+,
 \quad |h_r|\le H_rR.
 \label{eq:hbound}
\end{equation}
No derivative bound on the ODE matrix is required.

In $(K,L,N,J,w)$ order the conditioned difference data satisfy
\begin{equation}
 W_0=(h,0,0,h,w),\qquad W_{0,x}=(0,V_L,0,0,0),\qquad
 W_{0,r}=(h_r,0,0,h_r,w_r),
 \label{eq:initialvectors}
\end{equation}
where, with $k=p/(2a)$,
\begin{equation}
 w=k(h-2D_h/3),\qquad
 w_r=k_r(h-2D_h/3)+k(h_r-2D_{h,r}/3),\qquad
 V_L=h_r+bh.
 \label{eq:adapter}
\end{equation}
The mixed metric velocity and the derivative of $k$ are both retained.
Bounds following from \eqref{eq:caps} are
\begin{align}
 k_+&=\frac{1/7000}{2(111/1000)}=\frac1{1554},\nonumber\\
 k_{r,+}&=\ell\,131\left\{\frac{1/4000}{2(111/1000)}
              +\frac{(1/7000)(1/10000)}{2(111/1000)^2}\right\},\nonumber\\
 S&=k_+(H+2/3),\qquad
 S_r=k_{r,+}(H+2/3)+k_+(H_r+2/3),\qquad
 V=H_r+b_+H.
 \label{eq:componentcaps}
\end{align}
Thus $|w|\le SR$, $|w_r|\le S_rR$, and $|V_L|\le VR$.
For $\mathsf H=\operatorname{diag}(1,2,1,3,1)$, let
$I_0=\sup|W_0|_{\mathsf H}$, $I_r=\sup|W_{0,r}|_{\mathsf H}$,
$I_x=\sup|W_{0,x}|_{\mathsf H}$, with the data zero for $r\le1$.
The triangle inequalities give
\begin{equation}
 I_0\le(2H+S)R,\qquad I_r\le(2H_r+S_r)R,\qquad
 I_x\le2VR,\qquad
 N:=\max\{(I_0+TI_x)/3,I_r+I_x\}.
 \label{eq:N}
\end{equation}
All operations are rational. The resulting exact upper bound is
\begin{equation}
 \begin{gathered}
 N\le c_NR\le R/2000,\\
 c_N=
 \frac{798916710779559466175931203077}
      {1861847711577000000000000000000000}
 <0.000429098849392<\frac1{2000}.
 \end{gathered}
 \label{eq:Nrational}
\end{equation}
In particular every member of \eqref{eq:class} satisfies $N\le10^{30}$.
The certificate verifies 38 finite checks using standard-library rational
arithmetic, with zero new background or evolved-field oracle calls, PDE
solves, or bump queries. A separate rational implementation reconstructs
the bound from shared retained inputs without importing the producer
\cite{records}. These checks share the analytic and background premises.

\section{Transport to proper volume and expansion}\label{sec:transfer}

The shell observables use the intrinsic future normal-geodesic foliation
anchored at the common initial section:
\begin{equation}
 V_\eps(\tau)=\int_{S^3}\dd\mu_{h_\eps(\tau)},\qquad
 \Theta_\eps=\partial_\tau\log V_\eps,\qquad
 v_2=\frac{\partial_\eps^2V_\eps|_0}{V_0},\quad
 \Theta_2=\partial_\eps^2\Theta_\eps|_0.
 \label{eq:observables}
\end{equation}
The clock and anchor are transported under gauge changes. The background
has $A(\tau)=\rho_b\cosh(\tau/\rho_b)$, $V_0=2\pi^2A^3$, and
$\Theta_0=3\rho_b^{-1}\tanh(\tau/\rho_b)$.

The same first field means the same quadratic bulk, shell, clock, and
measure sources. They cancel in the difference of the second responses.
Retained linear uniqueness gives the complete homogeneous five-mean
difference $W$, including metric--scalar coupling and Robin feedback.
On the backward-closed cone $0\le x\le T$, $0\le r\le T-x$, the retained
coefficient estimate has
\begin{equation}
 q=\frac{1201}{3200000}<1,\qquad
 X(U)=\max\{\sup|U|_{\mathsf H}/3,
                   \sup|U_x|_{\mathsf H},\sup|U_r|_{\mathsf H}\}.
 \label{eq:resolvent}
\end{equation}
For the free reflected wave $B$ of \eqref{eq:initialvectors},
$X(B)\le N$ and $X(W-B)\le qN/(1-q)$, hence $X(W)\le N/(1-q)$.
The exact zero wedge $x+r\le1$ then implies, for $m(x)=W_w(x,0)$,
\begin{equation}
 |m(x)|\le (x-1)_+\frac{N}{1-q}.
 \label{eq:zero-wedge}
\end{equation}
Only the initial segment up to $T$ enters this target estimate.

Let $\nu=v_2^h-v_2^{\rm base}$ and
$\theta=\Theta_2^h-\Theta_2^{\rm base}$. The retained proper-time
difference equation and zero initial conditions are
\begin{equation}
 \nu_{\tau\tau}+4(A_\tau/A)\nu_\tau-2\nu/A^2=(F'_b/2)m,
 \qquad \nu(0)=\nu_\tau(0)=0,\qquad \theta=\nu_\tau.
 \label{eq:volumeode}
\end{equation}
There is no additional $1/2$ multiplying $\nu$ or $\theta$: the
observable coefficients are second parameter derivatives. At the target,
\begin{align}
 \nu(T)&=(\rho_b\ell)^2(F'_b/2)\int_1^T k(T,x)m(x)\,\dd x,\nonumber\\
 \theta(T)&=\rho_b\ell(F'_b/2)\int_1^T k_T(T,x)m(x)\,\dd x,\nonumber\\
 k(T,x)&=\frac{\sinh(2\ell(T-x))}{2\ell}
                   \frac{\cosh^2(\ell x)}{\cosh^2(\ell T)},\qquad
 k_T=\frac{\cosh^2(\ell x)\cosh(\ell T-2\ell x)}{\cosh^3(\ell T)}.
 \label{eq:kernels}
\end{align}
The exact inequalities $0\le k\le T-x$ and $0<k_T\le1$, together with
\eqref{eq:zero-wedge}, give
\begin{equation}
 \begin{gathered}
 |\nu(T)|\le C_vN,\qquad |\theta(T)|\le C_\Theta N,\\
 C_v=\frac{(131\ell)^2(1/1000)\Delta^3}{12(1-q)}
     =\frac{17161}{767711760000000000},\qquad
 C_\Theta=\frac{131\ell(1/1000)\Delta^2}{4(1-q)}
     =\frac{131}{63975980000}.
 \end{gathered}
 \label{eq:transport}
\end{equation}
This bounds the full coupled response, rather than a free-history sign.

\begin{theorem}[Conditional robustness of the specified class]\label{thm:main}
Under the retained premises, every admitted exact family with complete
increment $h\in\mathcal C_A$ has, at \eqref{eq:target},
\begin{equation}
 v_2\in[-7.6900\times10^{16},-4.1600\times10^{15}],\qquad
 \Theta_2\in[-1.1070\times10^{22},-2.3980\times10^{21}].
 \label{eq:mainbounds}
\end{equation}
For each such family there exists a possibly family-dependent
$\eps_h>0$ such that, for $0<|\eps|<\eps_h$ of either sign,
\begin{equation}
 0<V_\eps(\tau_*)<V_0(\tau_*),\qquad
 0<\Theta_\eps(\tau_*)<\Theta_0(\tau_*).
 \label{eq:qualitative}
\end{equation}
No positive parameter radius common to the whole class is asserted.
\end{theorem}

\begin{proof}
The retained base coefficient enclosures and the absolute difference
bounds at $N=10^{30}$ are
\begin{center}
\begin{tabular}{@{}lcc@{}}\toprule
Quantity & Base outward enclosure & Absolute increment cap\\\midrule
$v_2$ & $[-5.4520,-2.6521]\,10^{16}$ & $C_v10^{30}<2.23535\,10^{16}$\\
$\Theta_2$ & $[-9.0208,-4.4458]\,10^{21}$ & $C_\Theta10^{30}<2.04765\,10^{21}$\\
\bottomrule
\end{tabular}
\end{center}
The exact base rational intervals plus and minus the exact caps give
\eqref{eq:mainbounds} with outward rounding. Equivalently, $10^{30}$
is strictly below the retained sufficient threshold
$N_{\rm crit}>1.1864\times10^{30}$ for both coefficients to stay negative.
The first integrated variations vanish by the common quadrupole mean.
For each fixed admitted family, $C^2$ dependence gives
\begin{equation}
 V_\eps/V_0=1+\tfrac12\eps^2v_2+o(\eps^2),\qquad
 \Theta_\eps=\Theta_0+\tfrac12\eps^2\Theta_2+o(\eps^2).
 \label{eq:taylor}
\end{equation}
Since $V_0>0$ and $\Theta_0(\tau_*)>0$, strict negativity gives the
upper inequalities and continuity gives positivity after a further
family-dependent shrinkage. This proves \eqref{eq:qualitative}.
\end{proof}

The small-parameter conclusion is qualitative. The class bounds no higher
norms or nonlinear remainders uniformly, so the common coefficient
enclosures do not provide a uniform finite-$\eps$ regime. They concern
one registered proper time, not a signed time interval, contraction, or
the sign of volume acceleration.

\section{The earlier reversal and the scope of the result}

The earlier sign reversal uses the fixed admissible direction
\begin{equation}
 \begin{gathered}
 Q(r)=\exp\!\left(4-\frac1{(r-1)(2-r)}\right)\ (1<r<2),
 \qquad Q=0\text{ otherwise},\\
 \chi=\ell\rho_b^3Q,\qquad
 h_*=-\left(\rho_b/\rho\right)^3Q_r,
 \end{gathered}
 \label{eq:witness}
\end{equation}
scaled by the fixed finite $\lambda=-10^{21}$. Its two proper-response
coefficients are positive in that retained construction
\cite{preparation}. This does not contradict Theorem~\ref{thm:main}.
Independent exact rational target bounds give
\begin{equation}
 D_{h_*}(T)>3.7621\times10^{12},\qquad
 D_{h_*,r}(T)<-6.870\times10^{13}.
 \label{eq:witnessbounds}
\end{equation}
Consequently
\begin{equation}
 R_A(-10^{21}h_*)>3.7621\times10^{33}>2\times10^{33},\qquad
 |D_{-10^{21}h_*,r}(T)|>6.870\times10^{34}.
 \label{eq:excluded}
\end{equation}
It is outside the specified class already at one endpoint. This exclusion
uses retained target background boxes and an independent analytic Taylor
enclosure for the old $Q(T)$ and its derivatives; it is a check of an
existing witness, not a new direction search or oracle call.

The geometric ingredients have established antecedents. General first
and second variations of induced hypersurface geometry are treated by
Mars \cite[Secs.~3, 6--7]{Mars}; the cancellation in
Lemma~\ref{lem:area} is an application with a common first tangent.
Scalar--metric invariant combinations in warped scalar backgrounds are
standard \cite[Secs.~III--IV]{Giovannini}: with transverse conformal
coefficient $J=2\psi$, our $D=3(\psi-aw/p)$ has the same kinematic
cancellation. This does not identify the two dynamical problems.
The fact that a fixed scalar-defined window can give a gauge-invariant
averaging prescription does not select that window or a physical
hypersurface \cite[Sec.~3, Eq.~(3.23), and Sec.~4]{Gasperini}.
That work uses a timelike spacetime scalar gradient; here spatial scalar
levels lie inside an already supplied Cauchy section. Accordingly, the result
here is not a new gauge principle. It identifies the retained mean
freedom through a particular level-area profile, proves its restricted
reconstruction, and transports a stated profile radius to simultaneous
coefficient bounds in this model.

The supplied Cauchy section, continued level component, background label
calibration, canonical-sector restriction, and class radius are explicit
inputs. No physical process selecting them is given. The larger freedom
still admits the earlier opposite-sign preparation; the class theorem
states what follows after a further specification, not sign universality
from the first packet alone. No SI calibration, measured amplitude,
radiation epoch, or cosmological origin follows from these model-unit
coefficients. This focused comparison does not establish priority or
exhaustive literature novelty. The conditional mathematical result is
intended for author and specialist review.

\appendix
\section{Identification of the fixed seed and base}\label{app:base}

For clarity, the retained seed can be specified without choosing a
perturbed foliation. Let $Q$ be \eqref{eq:witness}, $B(z)=Q(z/\ell)$,
$p_b$ the background shell scalar derivative, and
\begin{gather}
 \alpha=(\rho/\rho_b)^{3/2}p/p_b,\qquad
 \varkappa=\alpha_z/\alpha=\rho(5a/2-\cV'/p),\nonumber\\
 u=\alpha B/\rho^2,\quad u_t=e^{-z}u,\quad s_1=u_t/\rho,\qquad
 v=3\alpha(B_z+\varkappa B)/(\rho^3p),\nonumber\\
 v_t=3\alpha e^{-z}[B_z+(\varkappa-1)B]/(\rho^3p),\qquad j=v_t/\rho.
 \label{eq:seed}
\end{gather}
In a background orthonormal spatial frame, with spatial density $\rho_0$,
the canonical first datum is
\begin{equation}
 \gamma_1=\operatorname{diag}(-4u,2u,2u,2u)Y,\quad
 \pi_1/\rho_0=\operatorname{diag}(-3s_1/2,0,0,0)Y,\quad
 \phi_1=vY,\quad\varpi_1/\rho_0=jY.
 \label{eq:canonicalseed}
\end{equation}
The mean Hamiltonian and radial momentum Hessians in convention
\eqref{eq:normalization} are
\begin{align}
 Q_H={}&-2u_y^2-24auu_y-2uu_{yy}
 -(16a_y+32a^2+8p^2/3)u^2-4u^2/\rho^2-s_1^2\nonumber\\
 &-(j^2+v_y^2+\cV''v^2)/6-(4/3)puv_y-(4/3)v^2/\rho^2,\nonumber\\
 Q_M={}&-2s_1u_y+j(v_y+4pu)/3.
 \label{eq:hessians}
\end{align}
Define $I=\int_{y_L}^{y_U}\rho^4Q_H\,\dd y$,
$J_B=\int_{y_L}^{y_U}\rho^2B\,\dd y>0$, $f=2IB/(9J_B)$, and
$\xi=-(\rho^4p^2)^{-1}\int_{y_L}^y\rho^4(Q_H+C[f])\,\dd s$, with
$C$ as in \eqref{eq:constraint}. The base second mean, in
$(\gamma,\pi,\phi,\varpi)$ order, is
\begin{equation}
 B_{\rm base}=\left(f\gamma_0+\Lie_{\xi\partial_y}\gamma_0,\,
 \rho_0\operatorname{diag}(0,-Q_M/(6a),-Q_M/(6a),-Q_M/(6a)),\,0,\,0\right).
 \label{eq:base}
\end{equation}
It satisfies the full mean second constraint with these Hessians.
The retained exact completion and evolution, rather than just the
initial Hessians, supply the base response intervals used in the theorem.
The added datum \eqref{eq:raw} leaves all canonical entries except the
metric mean unchanged.

\section{Reproducibility and shared inputs}

The accompanying capsule \cite{records} contains the geometric proof,
frozen design, rational producer, normal and optimized reports,
independent reconstruction, and the previous-witness exclusion check.
The rational producer authenticates the retained static annulus jets,
base proper response, and preparation-difference records by SHA-256
before evaluating the fixed gates. Normal and optimized runs agree in
their report bytes. The independent checker imports no producer.
Its agreement checks finite arithmetic under shared background and
analytic assumptions, not an independent observation or proof of the
Einstein--scalar PDE.

The scalar-level geometry and reconstruction require no new background
solution, evolved-field evaluation, numerical quadrature, or bump
evaluation. All new class bounds are rational consequences of
\eqref{eq:caps}, \eqref{eq:reconstruction}, and the retained coupled
estimate. The supplementary old-witness check constructs analytic target
exponential bounds and does not query a numerical bump oracle. Source,
displayed-bound audit, deterministic compilation record, and page-render
review accompany this manuscript.

\begin{thebibliography}{8}
\footnotesize
\setlength{\itemsep}{3pt}
\bibitem{Mars} M. Mars, \emph{First and Second Order Perturbations of
Hypersurfaces}, Class. Quantum Grav. \textbf{22} (2005) 3325--3348,
\href{https://arxiv.org/abs/gr-qc/0507005}{arXiv:gr-qc/0507005}.
\bibitem{Giovannini} M. Giovannini, \emph{Gauge-invariant fluctuations
of scalar branes}, Phys. Rev. D \textbf{64}, 064023 (2001),
\href{https://arxiv.org/abs/hep-th/0106041}{arXiv:hep-th/0106041}.
\bibitem{Gasperini} M. Gasperini, G. Marozzi, and G. Veneziano,
\emph{Gauge invariant averages for the cosmological backreaction},
JCAP \textbf{03} (2009) 011,
\href{https://arxiv.org/abs/0901.1303}{arXiv:0901.1303}.
\bibitem{preparation} R. Maldonado, \emph{Preparation dependence of
proper-time volume in Einstein--scalar shell perturbations}, working
manuscript (4 November 2026),
\href{https://github.com/maldonado-research/HDblast/tree/main/research/HDBLAST_SHELL_PREPARATION_DEPENDENCE_20261104}{preparation-difference capsule}.
\bibitem{records} R. Maldonado, \emph{HDBLAST conditional technical
records and reproducibility materials} (2026),
\href{https://github.com/maldonado-research/HDblast}{repository}.
The retained inputs include the localized constraint and development
records, the 1 November prepared-shell volume capsule, and the
4 November preparation-difference capsule. The present companion is
\href{https://github.com/maldonado-research/HDblast/tree/main/research/HDBLAST_INTRINSIC_PREPARATION_ROBUSTNESS_20261104}{the scalar-level preparation-robustness capsule}.
\end{thebibliography}
\end{document}
